\chapter{Validation and Debugging}
\label{ch:val}

Firmware that has only been compiled has not been validated. This chapter
reports the hardware campaign that was used both to find the defects of
\cref{ch:dev} and to confirm that the repairs hold on the bus. The
method is deliberately low-level: PCAN-View injects frames, the SR5E1
raises an FDCAN interrupt, TRACE32 steps through \texttt{DYNO2DWProtocol},
and PCAN-View records the response.

\section{Debug chain}
\label{sec:debug-chain}

\begin{figure}[ht]
\centering
\begin{tikzpicture}[
  font=\sffamily\footnotesize,
  box/.style={draw=NavyInk, rounded corners=1.5pt, align=center,
              inner sep=6pt, minimum height=11mm, minimum width=28mm},
  arr/.style={-{Latex}, thick, NavyInk}
]
  \node[box, fill=Paper] (pcan) {1. Frame injection\\PCAN-USB FD};
  \node[box, fill=Paper, right=12mm of pcan] (isr)
    {2. FDCAN IRQ\\Stellar-E};
  \node[box, fill=Paper, right=12mm of isr] (t32)
    {3. TRACE32 step\\DYNO2DWProtocol()};
  \node[box, fill=Paper, right=12mm of t32] (ack)
    {4. Bus response\\0x10 / 0x11};
  \draw[arr] (pcan) -- (isr);
  \draw[arr] (isr) -- (t32);
  \draw[arr] (t32) -- (ack);
\end{tikzpicture}
\caption{Closed-loop debug path from a synthetic CAN frame to the
acknowledgement observed on the same bus.}
\label{fig:debug-chain}
\end{figure}

The four steps of \cref{fig:debug-chain} were applied to every command in
\cref{tab:cmd-map}.

\paragraph{Frame injection.}
A PCAN-USB FD adapter, driven by PCAN-View~\cite{pcanview}, transmits a
frame with a chosen identifier, DLC and payload. \Cref{fig:pcan-tx} shows
the transmit list used during the campaign: \texttt{EXECUTE\_COMMAND}
(\texttt{0x33}), \texttt{SET\_REG} (\texttt{0x31}), \texttt{GET\_REG}
(\texttt{0x32}), the speed ramp, the current commands, and a handful of
diagnostic identifiers configured as receive filters.

\begin{figure}[ht]
\centering

\includegraphics[width=0.82\textwidth]{pcan_tx.png}
\caption{PCAN-View transmit list used for command injection. Each row is a
hand-composed CAN frame with identifier, DLC and payload bytes.}
\label{fig:pcan-tx}
\end{figure}

\paragraph{Interrupt.}
Arrival of a filtered identifier raises
\texttt{IRQ\_CAN1\_LINE0\_HANDLER} on the SR5E1. TRACE32 is programmed
with a breakpoint on the first instruction of
\texttt{DYNO2DWProtocol}, so the core halts with the payload already in
\texttt{RxArr[]}.

\paragraph{Step-by-step execution.}
From that breakpoint the engineer can watch the construction of
\texttt{msg[]}, the CRC, the call into \texttt{CFCP\_RX\_IRQ\_Handler},
and the Motor Control function that ultimately runs
(\texttt{UI\_SetReg}, \texttt{UI\_SetIdRef}, \texttt{MCI\_ExecFluxRamp},
\ldots). This is also how the SET\_REG discrepancy of
\cref{sec:discrepancy} was diagnosed: the first payload byte that the
specification called a length was being used as a register identifier,
and the TRACE32 locals window made that obvious in one step.

\paragraph{Response.}
On resume, the MCU publishes either identifier \texttt{0x10} with payload
\texttt{0xF0} (ACK) or \texttt{0xFF} (NACK), or identifier \texttt{0x11}
with a four-byte register value. PCAN-View records the frame; a mismatch
against the expected identifier or payload is a failed test.

\section{Physical layer and periodic diagnostics}
\label{sec:phys}

Before any command was trusted, the TX path was checked in isolation. With
the motor idle and no frames being injected, PCAN-View was configured to
log the diagnostic window \texttt{0x40}--\texttt{0x60}. The expected
behaviour, from \cref{sec:send-status}, is ten frames every
\SI{100}{\milli\second}.

\begin{figure}[ht]
\centering
\includegraphics[width=0.88\textwidth]{pcan_rx.png}
\caption{PCAN-View trace of MCU-to-dyno traffic. Identifier \texttt{0x10}
carries ACK/NACK; identifier \texttt{0x11} carries \texttt{REG\_VALUE};
the remaining identifiers are the periodic diagnostic stack.}
\label{fig:pcan-rx}
\end{figure}

The trace in \cref{fig:pcan-rx} confirms three properties at once.

\begin{enumerate}[leftmargin=1.6em]
  \item The physical layer, the bit-rate and the termination are correct:
        frames arrive without error frames and without ACK-slot errors.
  \item \texttt{Send\_Status} is running at the designed period. Counting
        the ten diagnostic identifiers between two occurrences of the
        sequence counter is a direct measurement of the \SI{100}{\milli\second}
        tick.
  \item The ACK/NACK identifier \texttt{0x10} and the \texttt{REG\_VALUE}
        identifier \texttt{0x11} are distinct, which is the observable
        counterpart of \cref{lst:ack}.
\end{enumerate}

A board that fails this test is not ready for command injection. Several
early bring-up failures---wrong bit-rate, missing 120~$\Omega$
termination, TX FIFO left in a non-idle state---were caught here, in
minutes, rather than inside TRACE32.

\section{Command-by-command tracking}
\label{sec:tracking}

For each command the same protocol was followed: inject, halt, step,
resume, compare. The notes below are the ones that changed the code.

\begin{itemize}[leftmargin=1.6em]
  \item \texttt{EXECUTE\_COMMAND 0x33}. The original start/stop complement
        was visible as a \texttt{STOP\_MOTOR} FCP frame being built for a
        CAN payload that had only \texttt{FAULT\_ACK} set. After the
        rewrite of \cref{sec:exec}, a payload with a single flag produces
        a single FCP frame with the matching \texttt{MC\_PROTOCOL\_CMD\_*}
        code.
  \item \texttt{SET\_RAMP 0x37}. Duration and final speed are unpacked
        little-endian; swapping the 16-bit duration bytes is a common
        DBC error and was caught here before the DBC was declared frozen.
  \item \texttt{SET\_CURRENT\_ID / IQ}. The 16-bit signed reconstruction
        was checked with the values \texttt{0x00E8} ($+232$) and
        \texttt{0xFF18} ($-232$) to confirm sign extension.
  \item \texttt{SET\_REG 0x31}. See the next section; this is the case
        that disagreed with the written protocol.
  \item PI round-trip. A \texttt{SET\_REG} of \texttt{SPEED\_KP} with
        numerator $N$ and divisor $n$, followed by a \texttt{GET\_REG} of
        the same register, returned the packed word
        $(n \ll 16) \mid N$ on identifier \texttt{0x11}.
\end{itemize}

\section{Case study: SET\_REG versus the dishwasher specification}
\label{sec:discrepancy}

The Motor Control protocol document~\cite{stmcprotocol} describes
\texttt{SET\_REG} as a length-prefixed frame: byte~0 is the payload
length, byte~1 is the register identifier, the following bytes are the
value. The SR5E1 library, as written and as executed, does something else.
Byte~0 of the CAN payload \emph{is} the register identifier. There is no
length byte on the wire; the FCP length is inserted by DYNO2DW when it
builds \texttt{msg[]}.

\begin{figure}[ht]
\centering
\begin{minipage}[t]{0.48\textwidth}
  \centering
  \includegraphics[width=\textwidth]{setreg_pcan.png}
  \captionof{figure}{SET\_REG frame as transmitted from PCAN-View.
  Byte~0 of the payload is already the target register identifier.}
  \label{fig:setreg-pcan}
\end{minipage}\hfill
\begin{minipage}[t]{0.48\textwidth}
  \centering
  \includegraphics[width=\textwidth]{dw_spec.png}
  \captionof{figure}{SET\_REG layout in the dishwasher protocol
  specification, with a length byte in front of the register identifier.}
  \label{fig:dw-spec}
\end{minipage}
\end{figure}

TRACE32 made the disagreement uncontroversial. Halting in the
\texttt{SET\_REG} case showed \texttt{RxArr[0]} being copied into the
register-id argument of \texttt{UI\_SetReg}, not into \texttt{msg[1]}
as a length. A frame built according to the specification---length first---
wrote the wrong register and returned a NACK or, worse, a silent write
to an unintended PI gain.

Two equally defensible repairs were available: change the firmware to
match the document, or document the firmware and change the DBC. The
second option was taken, for three reasons.

\begin{enumerate}[leftmargin=1.6em]
  \item The library was already deployed on boards that the dyno spoke to.
        Changing the on-wire layout would have broken those boards and
        every recorded trace.
  \item DYNO2DW exists specifically to absorb this kind of difference: the
        FCP frame, with its length byte and CRC, is an internal MCU
        affair. The CAN payload can be denser.
  \item The DBC is the contract the GUI and the dyno software compile. Once
        it says ``byte~0 is \texttt{REGISTER\_ID}'', both consumers are
        right together.
\end{enumerate}

The episode is recorded at this length because it is the standard failure
mode of protocol work: two artefacts that each look consistent, and a bus
that cannot tell them apart until a debugger sits on both sides of the
ISR.

\section{What was considered passing}
\label{sec:pass-criteria}

A command was marked validated when all of the following held on hardware.

\begin{itemize}[leftmargin=1.6em]
  \item The TRACE32 locals showed the intended FCP code, length and
        payload, and a CRC that \texttt{SelfCRCCalc} accepted.
  \item The Motor Control function associated with that code executed
        (verified by a breakpoint on the function itself, not only on
        DYNO2DW).
  \item The bus showed identifier \texttt{0x10} with \texttt{0xF0} for
        writes, or identifier \texttt{0x11} with the expected four-byte
        value for reads.
  \item A second, malformed variant of the same command (truncated DLC,
        reserved flag set, out-of-range register id) produced
        \texttt{0xFF} and did not change machine state.
\end{itemize}

The GUI of \cref{ch:gui} was not used as an oracle during this campaign.
It was connected only after the DBC matched the traces, so that a widget
bug could not be mistaken for a firmware bug.
